% Chapter 5: Generalized Perspective Functions
\chapter{Generalized Perspective Functions}

This chapter establishes the joint convexity and joint concavity of the generalized
perspective function $(f \triangle g)$ defined in Chapter~1
(see Definition~\ref{def:gen_perspective}).
The two main theorems --- Theorem~\ref{thm:perspective_joint_convex} and
Theorem~\ref{thm:perspective_joint_concave} --- are abstract C$^*$-algebra results
that simultaneously generalize Lieb's concavity theorem and Ando's convexity theorem.
The power-mean corollaries that close the chapter are the specializations used in
Chapter~6.

\section{Joint Convexity}

\begin{theorem}[Joint convexity of the generalized perspective function]
  \label{thm:perspective_joint_convex}
  \lean{PerspectiveJointConvex}
  \uses{def:operator_convex, def:operator_concave, def:gen_perspective}
  Let $\mathcal{B}$ be an ordered unital C$^*$-algebra, and let
  $f, g : \mathbb{R} \to \mathbb{R}$ be continuous on $[0,\infty)$.
  Assume:
  \begin{itemize}
    \item $f$ is operator convex on $[0,\infty)$ with $f(0) \leq 0$,
    \item $g$ is operator concave on $[0,\infty)$ with $g(x) > 0$ for all $x > 0$.
  \end{itemize}
  Then the map $(L, R) \mapsto (f \triangle g)(L, R)$ is jointly convex on the set
  \[
    \{ (L, R) \in \mathcal{B} \times \mathcal{B} \mid 0 \leq L,\ R \text{ strictly positive} \}.
  \]
  \begin{proof}
    \uses{thm:jensen_sub, thm:loewner_concave}
    Fix pairs $(L_1, R_1)$ and $(L_2, R_2)$ in the domain and weights $a, b \geq 0$
    with $a + b = 1$.
    Let $R := a R_1 + b R_2$ and write $G := \operatorname{cfc}(g)$ for the operator
    $g$ applied via the continuous functional calculus.
    Define bounded operators
    \[
      T_1 := \bigl(a\, G(R_1)\bigr)^{1/2} G(R)^{-1/2},
      \qquad
      T_2 := \bigl(b\, G(R_2)\bigr)^{1/2} G(R)^{-1/2}.
    \]
    Because $g$ is operator concave, the operator concavity inequality gives
    $a\, G(R_1) + b\, G(R_2) \leq G(R)$, and therefore
    \[
      T_1^* T_1 + T_2^* T_2
      = G(R)^{-1/2}(a\, G(R_1) + b\, G(R_2))G(R)^{-1/2}
      \leq \mathrm{id}.
    \]
    That is, the pair $(T_1, T_2)$ is \emph{sub-unital} in the sense required by
    the sub-unital Jensen inequality (Theorem~\ref{thm:jensen_sub}).
    Applying Theorem~\ref{thm:jensen_sub} to the operator-convex function $f$
    with arguments $G(R)^{-1/2} L_i\, G(R)^{-1/2}$ and operators $T_i$, then
    conjugating by $G(R)^{1/2}$ and collecting terms, yields
    \[
      a\,(f \triangle g)(L_1, R_1) + b\,(f \triangle g)(L_2, R_2)
      \;\geq\;
      (f \triangle g)(aL_1 + bL_2,\, R).
    \]
    (The condition $f(0) \leq 0$ absorbs the sub-unital defect term in Jensen's
    inequality.)
    This is precisely the joint convexity inequality.
  \end{proof}
\end{theorem}

\section{Joint Concavity}

\begin{theorem}[Joint concavity of the generalized perspective function]
  \label{thm:perspective_joint_concave}
  \lean{PerspectiveJointConcave}
  \uses{def:operator_concave, def:gen_perspective}
  Let $\mathcal{B}$ be an ordered unital C$^*$-algebra, and let
  $f, g : \mathbb{R} \to \mathbb{R}$ be continuous on $[0,\infty)$.
  Assume:
  \begin{itemize}
    \item $f$ is operator concave on $[0,\infty)$ with $f(0) \geq 0$,
    \item $g$ is operator concave on $[0,\infty)$ with $g(x) > 0$ for all $x > 0$.
  \end{itemize}
  Then the map $(L, R) \mapsto (f \triangle g)(L, R)$ is jointly concave on the set
  \[
    \{ (L, R) \in \mathcal{B} \times \mathcal{B} \mid 0 \leq L,\ R \text{ strictly positive} \}.
  \]
  \begin{proof}
    \uses{thm:perspective_joint_convex}
    Observe that $(-f) \triangle g = -(f \triangle g)$ and that $-f$ is operator
    convex with $(-f)(0) \leq 0$; apply Theorem~\ref{thm:perspective_joint_convex}
    to $-f$ and negate to obtain joint concavity of $f \triangle g$.
  \end{proof}
\end{theorem}

\section{Power Mean Monotonicity}

The \emph{$(\alpha,\beta)$-operator power mean} is the special case of the generalized
perspective function with $f(t) = t^\alpha$ and $g(t) = t^\beta$.
By L\"{o}wner's theorem (Chapter~4), the exponent ranges determine operator
convexity or concavity of these power functions, and the following corollaries
follow at once.

\begin{theorem}[Joint concavity of the operator power mean]
  \label{thm:power_mean_concave}
  \lean{PowerMeanJointlyConcave}
  \uses{def:gen_perspective}
  For $0 < \alpha \leq 1$ and $0 < \beta \leq 1$, the map
  $(L, R) \mapsto ((\cdot)^\alpha \triangle (\cdot)^\beta)(L, R)$
  is jointly concave on
  $\{ (L, R) \in \mathcal{B} \times \mathcal{B} \mid 0 \leq L,\ R \text{ strictly positive} \}$.
  \begin{proof}
    \uses{thm:perspective_joint_concave, thm:loewner_concave}
    Both $t \mapsto t^\alpha$ and $t \mapsto t^\beta$ are operator concave on
    $[0,\infty)$ when the exponents lie in $(0,1]$ by L\"{o}wner's theorem
    (Theorem~\ref{thm:loewner_concave}); apply Theorem~\ref{thm:perspective_joint_concave}.
  \end{proof}
\end{theorem}

\begin{theorem}[Joint convexity of the operator power mean]
  \label{thm:power_mean_convex}
  \lean{PowerMeanJointlyConvex}
  \uses{def:gen_perspective}
  For $1 \leq \alpha \leq 2$ and $0 < \beta \leq 1$, the map
  $(L, R) \mapsto ((\cdot)^\alpha \triangle (\cdot)^\beta)(L, R)$
  is jointly convex on
  $\{ (L, R) \in \mathcal{B} \times \mathcal{B} \mid 0 \leq L,\ R \text{ strictly positive} \}$.
  \begin{proof}
    \uses{thm:perspective_joint_convex, thm:loewner_convex, thm:loewner_concave}
    For $\alpha \in [1,2]$ the function $t \mapsto t^\alpha$ is operator convex on
    $[0,\infty)$ with $0^\alpha = 0 \leq 0$ (Theorem~\ref{thm:loewner_convex}),
    and for $\beta \in (0,1]$ the function $t \mapsto t^\beta$ is operator concave
    (Theorem~\ref{thm:loewner_concave}); apply Theorem~\ref{thm:perspective_joint_convex}.
  \end{proof}
\end{theorem}
